\chapter{Aspectos éticos y protección de datos}
\label{chap:gestion_etica}

\section{Aspectos éticos}

El trabajo tiene carácter retrospectivo y metodológico. Su finalidad es comparar enfoques de aprendizaje con pocos datos etiquetados sobre imágenes de ECG, no automatizar decisiones asistenciales ni sustituir la revisión clínica. Cualquier predicción debe interpretarse como una señal experimental de priorización, nunca como diagnóstico, descarte de enfermedad o recomendación terapéutica.

La principal precaución ética es evitar conclusiones clínicas que el diseño no puede sostener. El síndrome de Brugada requiere valoración especializada, contexto clínico y, en pacientes seleccionados, pruebas complementarias. Por tanto, las métricas obtenidas en esta memoria solo describen el comportamiento de los modelos dentro de los datos y protocolos empleados.

Esta cautela se aplica de forma especial al enfoque VLM/ICL. Que un modelo modifique su decisión al recibir ejemplos no equivale a competencia clínica autónoma. En este trabajo, la mejora frente al zero-shot se interpreta como señal metodológica de adaptación al contexto, no como autorización para usar el sistema sin revisión profesional.

La misma cautela se aplica a la adaptación de dominio de la CNN. Una mejora modesta al usar los datos cedidos por el HUCA sin etiquetas morfológicas no constituye validación clínica. Solo indica que ciertas alineaciones de representación pueden reducir parcialmente la brecha entre simulador y datos reales.

También debe considerarse el posible sesgo de selección. El conjunto real procede de un recurso concreto y el tamaño de positivos es limitado. La ausencia de análisis por subgrupos no implica ausencia de sesgo, sino falta de tamaño suficiente para estudiarlo con garantías. Una extensión clínica exigiría validación externa, evaluación prospectiva y supervisión profesional.

\section{Protección de datos}

Los datos de salud son una categoría especial según el artículo 9 del Reglamento General de Protección de Datos \cite{gdpr}. Su tratamiento exige una base jurídica adecuada, minimización de datos, limitación de finalidad, confidencialidad y medidas de seguridad proporcionadas al riesgo. Para investigación científica, el artículo 89 exige salvaguardias técnicas y organizativas.

En España, la Ley Orgánica 3/2018 completa este marco y regula garantías para tratamientos con fines de investigación, especialmente cuando se trabaja con datos seudonimizados \cite{lopdgdd}. La Ley 14/2007 de Investigación biomédica añade principios de confidencialidad, control ético y protección de las personas participantes \cite{ley142007}.

El recurso Brugada HUCA se publica sin información identificativa y declara anonimización mediante Safe Harbor \cite{costacortez2026brugada}. Aun así, las imágenes derivadas de ECG deben tratarse con prudencia: solo deben conservarse los campos necesarios para el análisis, evitar identificadores directos en materiales de difusión y mantener los datos en entornos controlados. Si se incorporasen datos adicionales no anónimos, habría que revisar la base jurídica y realizar una evaluación de impacto antes del tratamiento.
